% rutowski.tex

\section{Radiating argon shock layer with thermochemical nonequilibrium}
\label{sec:}
\index{radiation transport model!example of use}
\index{finite-rate chemistry!example of use}
\index{energy exchange!example of use}
\index{gas model!two temperature!example of use}
%

\subsection{Introduction}

Rutowski \textit{et al.}~\cite{RB64} measured the total and radiative heat fluxes at the stagnation point of a 1\,inch diameter hemisphere immersed in a freestream flow of shock heated argon.
A schematic of the experimental setup is depicted in Figure~\ref{fig:schematic}.
The hemispherical model was placed in the test section of a 3 inch diameter stainless steel shock tube at the Lockheed Research Laboratories~\cite{MBR+62}.  
Incident shock waves with velocities up to 4.3\,km/s ($M_s = 13.2$) were driven through the argon test gas at an initial pressure of 10\,Torr.
Total heat transfer at the stagnation point was measured with a surface mounted calorimetric gauge, in which a thin strip of polished platinum is exposed to the flow and the heat transfer determined from change in resistivity.
Radiative heat transfer was measured with a similar gauge mounted behind a sapphire window that allowed transmission in the wavelength range $180 \leq \lambda \leq 6000$\,nm.
The platinum strip was determined to have a weighted average absorptivity of 0.4, however some experiments were also performed with a thin layer of camphor lampblack to give an absorptivity of 1.0.
The error in the total heat transfer and radiative measurements were estimated to be $\pm 5$\,\% and $\pm 15$\,\%, respectively. 
 

\begin{figure}
 \centering
 \includegraphics[scale=0.5]{../2D/Rutowski-hemisphere/figures/schematic.pdf}
 \caption{Schematic diagram of hemispherical model immersed in shock heat flow (adapted from Rutowski \textit{et al.}~\cite{RB64}).}
  \label{fig:schematic}
\end{figure}

In the present work a shock Mach number of 12.7 is considered.
The computational domain and boundary conditions are illustrated in Figure~\ref{fig:domain}.
The stainless steel surface is modelled as a fixed temperature, fully catalytic wall at 300\,K.
% chemical reactions
% energy exchange
% transport model


\begin{figure}
\centering
 \subfloat[Computation domain] {
  \includegraphics[width=0.49\linewidth]{../2D/Rutowski-hemisphere/figures/domain.pdf}
 }
 \subfloat[Computation grid] {
  \includegraphics[width=0.28\linewidth]{../2D/Rutowski-hemisphere/figures/grid.png} \vspace{1cm}
 }
 \caption{Computational domain and grid for Navier--Stokes simulations of the Rutowski and Bershader~\cite{RB64} experiments.}
  \label{fig:domain}
\end{figure}

% plot of results



\subsection{Eilmer3 input scripts (.py)}

\subsubsection{Part 1 -- inviscid flow}

\topbar
\lstinputlisting[language={}]{../2D/Rutowski-hemisphere/Ms_12.70/part1-inviscid/hemisphere.py}
\bottombar

\subsubsection{Part 2 -- viscous flow}

\topbar
\lstinputlisting[language={}]{../2D/Rutowski-hemisphere/Ms_12.70/part2-viscous/hemisphere.py}
\bottombar

\subsubsection{Part 3 -- viscous flow with radiation coupling}

\topbar
\lstinputlisting[language={}]{../2D/Rutowski-hemisphere/Ms_12.70/part3-viscous-with-radiation/hemisphere.py}
\bottombar

\subsection{Chemical reaction script (.lua)}

\topbar
\lstinputlisting[language={}]{../2D/Rutowski-hemisphere/kinetic-models/Ar-2T-chemical-reactions.lua}
\bottombar

\subsection{Thermal energy exchange script (.lua)}

\topbar
\lstinputlisting[language={}]{../2D/Rutowski-hemisphere/kinetic-models/Ar-2T-energy-exchange.lua}
\bottombar

\subsection{Radiation model script (.py)}

\topbar
\lstinputlisting[language={}]{../2D/Rutowski-hemisphere/Ms_12.70/part3-viscous-with-radiation/Ar-nonequilibrium-radiation.py}
\bottombar

\subsection{Tcl run script}
\topbar
\lstinputlisting[language={}]{../2D/Rutowski-hemisphere/Ms_12.70/Rutowski-long.test}
\bottombar

\subsection{Error checking script (.py)}
\topbar
\lstinputlisting[language={}]{../2D/Rutowski-hemisphere/Ms_12.70/compute_qrad_error.py}
\bottombar

\subsection{Notes}
